%% ============================================================
%%  Orleu Bachelor Thesis — Full LaTeX Source
%%  Astana IT University, Department of Software Engineering
%%  Authors: Nurmukhammed Kanafin, Adilet Serikbay, Baglan Tolegenov
%%  Supervisor: Aigerim Kydyrbekova
%% ============================================================

\documentclass[12pt, a4paper]{report}

%% ---------- Packages ----------
\usepackage[T1]{fontenc}
\usepackage[utf8]{inputenc}
\usepackage[english]{babel}
\usepackage{lmodern}
\usepackage{microtype}
\usepackage{geometry}
\usepackage{setspace}
\usepackage{parskip}
\usepackage{indentfirst}
\usepackage{titlesec}
\usepackage{titletoc}
\usepackage{fancyhdr}
\usepackage{emptypage}
\usepackage{graphicx}
\usepackage{float}
\usepackage{caption}
\usepackage{subcaption}
\usepackage{booktabs}
\usepackage{longtable}
\usepackage{array}
\usepackage{tabularx}
\usepackage{multirow}
\usepackage{xcolor}
\usepackage{hyperref}
\usepackage{url}
\usepackage{natbib}
\usepackage{amsmath}
\usepackage{enumitem}
\usepackage{pdfpages}

%% ---------- Page Geometry ----------
\geometry{
  a4paper,
  top=25mm,
  bottom=25mm,
  left=30mm,
  right=20mm
}

%% ---------- Line Spacing ----------
\onehalfspacing

%% ---------- Paragraph Indent ----------
\setlength{\parindent}{1.25cm}
\setlength{\parskip}{6pt}

%% ---------- Colors ----------
\definecolor{crimson}{RGB}{200, 52, 58}
\definecolor{darkgray}{RGB}{50, 50, 50}

%% ---------- Hyperref Setup ----------
\hypersetup{
  colorlinks = true,
  linkcolor  = darkgray,
  citecolor  = darkgray,
  urlcolor   = crimson,
  pdftitle   = {Development of a Gamified Fitness Application for Tracking Workouts, Nutrition, and Motivation through Game Mechanics},
  pdfauthor  = {Nurmukhammed Kanafin, Adilet Serikbay, Baglan Tolegenov}
}

%% ---------- Chapter / Section Titles ----------
\titleformat{\chapter}[display]
  {\normalfont\Large\bfseries\centering}
  {\chaptertitlename\ \thechapter}{12pt}{\Large}
\titlespacing*{\chapter}{0pt}{0pt}{20pt}

\titleformat{\section}
  {\normalfont\large\bfseries}{\thesection}{1em}{}
\titleformat{\subsection}
  {\normalfont\normalsize\bfseries}{\thesubsection}{1em}{}
\titleformat{\subsubsection}
  {\normalfont\normalsize\itshape}{\thesubsubsection}{1em}{}

%% ---------- Header / Footer ----------
\pagestyle{fancy}
\fancyhf{}
\fancyhead[L]{\small\nouppercase{\leftmark}}
\fancyhead[R]{\small\thepage}
\fancyfoot{}
\renewcommand{\headrulewidth}{0.4pt}

\fancypagestyle{plain}{
  \fancyhf{}
  \fancyfoot[C]{\small\thepage}
  \renewcommand{\headrulewidth}{0pt}
}

%% ---------- Table / Figure Captions ----------
\captionsetup{font=small, labelfont=bf, skip=6pt}

%% ---------- Enumitem ----------
\setlist[enumerate]{leftmargin=*, nosep}
\setlist[itemize]{leftmargin=*, nosep}

%% ============================================================
\begin{document}

%% ============================================================
%%  TITLE PAGE
%% ============================================================
\begin{titlepage}
  \centering
  \vspace*{10mm}
  {\Large\bfseries ASTANA IT UNIVERSITY\par}
  \vspace{4mm}
  {\large Department of Software Engineering\par}
  \vspace{8mm}
  {\large Bachelor Thesis\par}
  \vspace{20mm}
  {\LARGE\bfseries Development of a Gamified Fitness Application\\[6pt]
   for Tracking Workouts, Nutrition, and Motivation\\[6pt]
   through Game Mechanics\par}
  \vspace{20mm}
  \begin{flushleft}
    {\large\textbf{Authors:}}\par
    \vspace{4mm}
    {\large Nurmukhammed Kanafin\par}
    {\large Adilet Serikbay\par}
    {\large Baglan Tolegenov\par}
    \vspace{10mm}
    {\large\textbf{Program:} Software Engineering\par}
    \vspace{10mm}
    {\large\textbf{Supervisor:} Aigerim Kydyrbekova\par}
    {\large Senior Lecturer\par}
  \end{flushleft}
  \vfill
  {\large Astana, 2026\par}
\end{titlepage}

%% ============================================================
%%  FRONT MATTER
%% ============================================================
\pagenumbering{roman}
\setcounter{page}{2}

%% ---------- Abstract ----------
\chapter*{Abstract}
\addcontentsline{toc}{chapter}{Abstract}

This thesis presents the design and implementation of Orleu, a mobile fitness application that combines adaptive gamification with machine learning-based progress prediction. The central problem addressed is the well-documented inability of current fitness applications to sustain long-term user engagement: most applications rely on static reward mechanisms that fail to respond to how individual users actually progress over time. Orleu addresses this gap by using a LightGBM classification model to detect whether a user's training trajectory is improving, plateauing, or declining, and by dynamically adjusting mission difficulty, narrative content, and motivational feedback in response to those predictions.

The system incorporates workout logging, lightweight nutritional context tracking, and a gamified campaign structure in which users progress through narrative chapters and choose weekly missions. Nutritional consistency, measured as the proportion of days in a rolling fourteen-day window for which a nutrition entry was recorded, is included as one of five features in the ML pipeline, reflecting evidence that training adaptations are influenced by protein intake and dietary adherence. Predictions drive a three-level adaptation mechanism that scales mission targets between 75 and 115~percent of baseline depending on the detected trend. An AI coach component communicates the rationale for adaptations in plain language, addressing transparency concerns raised in the adaptive systems literature.

The system is built on React Native with Expo SDK~54 for the mobile client, FastAPI and PostgreSQL for the backend, and LightGBM with SHAP for the prediction layer. Evaluation is conducted through a structured 43-question questionnaire survey (\(N = 51\)) recruiting participants from the Astana IT University fitness community. The instrument captures training habits, motivational challenges, preferences for adaptive difficulty adjustment, and perceived value of AI-generated feedback, and is analysed using descriptive statistics with Likert-scale means and preference rankings mapped directly to system design decisions. Expected findings and design implications are discussed in relation to self-determination theory and flow theory.

\vspace{8pt}
\noindent\textbf{Keywords:} adaptive gamification, mobile health, machine learning, fitness tracking, LightGBM, self-determination theory, user motivation, progress prediction

%% ---------- Table of Contents ----------
\tableofcontents
\listoffigures
\listoftables

%% ============================================================
%%  MAIN MATTER
%% ============================================================
\cleardoublepage
\pagenumbering{arabic}

%% ============================================================
\chapter{Introduction}
%% ============================================================

\section{Background and Motivation}

Mobile fitness applications have become one of the most widely downloaded categories in app stores globally, yet their record at producing lasting behavior change remains poor. Surveys of wearable device usage consistently find that engagement drops sharply after the first few weeks of use, with a large proportion of users abandoning tracking applications within six months of download \citep{ledger2014}. The disconnect between initial adoption and sustained engagement points to a structural limitation in how most fitness software is designed: applications are optimized for acquisition rather than retention, and their motivational mechanisms are built once and then applied identically to every user regardless of fitness level, progress trajectory, or evolving preferences.

The most common motivational layer in commercial fitness applications is gamification, broadly defined as the use of game design elements in non-game contexts \citep{deterding2011}. Points, badges, streaks, and leaderboards appear across virtually every major fitness platform. These elements are not without effect. Research has shown that gamification can increase initial engagement and self-reported motivation in short-term interventions \citep{johnson2016}. The problem is that this effect does not hold over time. Users habituate to fixed reward schedules, and once the novelty of earning badges fades, the mechanics cease to function as motivators. A user who has trained consistently for six months faces the same badge thresholds and the same challenge difficulty as a user in their first week. The system has no mechanism to recognise that their situation has changed.

This limitation is not merely a design oversight. It reflects a deeper mismatch between static systems and the dynamic nature of human skill development. Flow theory, developed by Csikszentmihalyi and extended to digital contexts, predicts that engagement is highest when perceived challenge closely matches perceived ability \citep{nakamura2014}. When challenge falls below ability, boredom follows. When challenge exceeds ability, anxiety results. Static gamification is structurally incapable of maintaining this balance because it does not track how ability changes over time. Self-determination theory adds a parallel concern: motivational systems that feel controlling rather than supportive undermine the intrinsic motivation that sustains long-term behavior change \citep{ryan2000}. Points and badges that are awarded on fixed schedules, without regard for whether they represent a meaningful accomplishment for this particular user at this particular moment, tend to be perceived as external rewards rather than genuine recognition of progress.

The proposed solution is an adaptive system that monitors actual training behavior, detects changes in progress trajectory using machine learning, and adjusts gamification parameters in response. This is the core idea behind Orleu. Rather than treating all users identically, the system classifies each user's recent training data as improving, plateauing, or declining, and uses that classification to scale mission difficulty, select narrative content, and calibrate feedback tone. The adaptation is not prescriptive; the system does not tell users what to do. It responds to what users have already done and adjusts the motivational environment accordingly.

\section{The Role of Nutritional Context}

Workout logs in isolation offer an incomplete picture of an athlete's training status. Sport science has long established that protein intake and caloric availability are primary determinants of resistance training adaptation \citep{morton2018, helms2014}, meaning that a plateau in recorded performance could reflect genuine stagnation or simply inadequate dietary support during a period of energy restriction. These two conditions call for markedly different responses, yet they are indistinguishable when only exercise data is available.

Orleu addresses this gap through a dedicated nutrition tracking module. Users log individual meals across four categories of breakfast, lunch, dinner, and snacks by searching a food database and specifying quantities. Each entry stores pre-computed values for calories, protein, carbohydrates, and fat. Personalized daily targets can be configured through a goals interface, and the application renders progress against those targets both on a per-day basis and as aggregated trends over a rolling weekly view. The food database supports full-text fuzzy search via a GIN-indexed name field, reducing friction when items cannot be recalled with exact spelling. Users may also create and retain custom entries in a personal library, accommodating foods absent from the shared database.

Within the machine learning pipeline, nutritional consistency serves as one of five features supplied to the LightGBM classifier. This feature is defined as the proportion of days in a fourteen-day rolling window for which at least one nutrition log was recorded. The decision to use consistency as a proxy rather than raw macronutrient totals reflects the practical observation that the behavioral act of tracking---independent of the precise values recorded---is associated with greater dietary awareness and adherence. This makes the feature meaningful even when users log incompletely or imprecisely.

\section{Project Objectives}

The objectives of this project are:

\begin{enumerate}
  \item To design and implement a mobile fitness application that logs workout data including exercises, sets, repetitions, and load, and captures lightweight nutritional context for use in predictive modeling.
  \item To develop a LightGBM classification model that predicts user progress trend from structured training and nutritional features, and to validate that model using both synthetic benchmark data and real user data collected during the evaluation period.
  \item To implement an adaptive gamification engine that adjusts mission difficulty, narrative content, and feedback tone based on ML predictions, in a manner consistent with the challenge-skill balance principle from flow theory and the autonomy-support principles from self-determination theory.
  \item To evaluate the proposed system with real users from the university fitness community through a structured survey (\(N = 51\)) measuring training habits, motivational challenges, preferences for adaptive gamification features, and perceived value of AI-generated feedback.
  \item To document design decisions, architectural constraints, and evaluation findings in a form that contributes to the literature on adaptive gamification in digital health.
\end{enumerate}

\section{Scope and Limitations}

The scope of this project is a gym-based strength training application for adult users without clinical health conditions. The application targets recreational and semi-serious weightlifters rather than competitive athletes or rehabilitation patients. Nutritional tracking is limited to daily caloric and protein values; micronutrients, meal timing, and hydration are outside scope.

The ML prediction model requires a minimum of fourteen days of engagement and at least three logged workouts before generating predictions. Users who have not met this threshold receive static gamification without adaptive adjustment. This cold start period is a known limitation of the system and is discussed in Chapter~\ref{ch:limitations}.

The evaluation study is limited by the constraints of a bachelor's thesis timeline. The survey-based evaluation (\(N = 51\)) captures user preferences and expectations rather than longitudinal behaviour change. This is not large enough to support causal claims about behavior change. The study is designed to generate preliminary evidence of engagement effects and to identify usability issues for future development.

\section{Thesis Structure}

Chapter~\ref{ch:litreview} reviews literature on gamification in fitness applications, adaptive persuasive systems, machine learning for mobile health, nutritional context in performance modeling, and motivational psychology. Chapter~\ref{ch:methodology} describes the research methodology and evaluation design. Chapter~\ref{ch:implementation} presents the system architecture and implementation details across the ML pipeline, gamification engine, mobile application, and backend. Chapter~\ref{ch:results} reports ML model validation results and the user study design with expected findings. Chapter~\ref{ch:conclusion} concludes the thesis. Chapter~\ref{ch:limitations} discusses limitations and directions for future work.

%% ============================================================
\chapter{Literature Review}
\label{ch:litreview}
%% ============================================================

\section{Overview}

This chapter reviews literature across five domains relevant to the design of Orleu: gamification in fitness applications, adaptive persuasive technology, machine learning in mobile health, nutritional context in training performance, and motivational psychology. The search covered ACM Digital Library, IEEE Xplore, Springer Link, PubMed, and Google Scholar, prioritizing peer-reviewed journal articles and conference proceedings published between 2012 and 2024. Studies focusing exclusively on clinical populations or serious games rather than gamification were excluded. The chapter concludes by identifying the research gap that Orleu addresses and positioning the proposed system within the existing body of knowledge.

\section{Gamification in Fitness Applications}

\subsection{Defining Gamification in Health Contexts}

\citet{hamari2014} define gamification as the implementation of game design elements in non-game contexts with the intent of influencing user behavior. This distinguishes gamification from serious games, which are complete game experiences designed for non-entertainment purposes. In fitness applications, the core activity remains workout logging and exercise performance; game elements are layered on top to increase engagement rather than replacing the primary task. \citet{cugelman2013} notes that this distinction has practical consequences for digital health developers, as gamification can be integrated into existing tracking systems without requiring users to learn entirely new interaction paradigms.

The most common gamification elements in fitness applications are point systems, badges, leaderboards, and progress bars, often abbreviated as PBL mechanics. These elements appear across virtually every major commercial fitness platform. \citet{deterding2015} argues that the prevalence of PBL mechanics reflects the relative ease of implementing them rather than their demonstrated effectiveness, and proposes a more principled approach grounded in the concept of intrinsic skill atoms that match game elements to specific psychological needs.

\subsection{Empirical Evidence of Effectiveness}

\citet{johnson2016} conducted a systematic review of 19 empirical gamification studies in health and wellbeing contexts and found positive effects on engagement and motivation in the majority of studies, particularly in short-term interventions lasting less than eight weeks. \citet{hamari2014} reviewed 24 studies and reached a similar conclusion: gamification tends to produce measurable positive effects, but the effects are highly context-dependent and diminish over time. \citet{sardi2017} extended this analysis to e-health specifically and found that while 83~percent of reviewed applications reported positive engagement metrics, only 29~percent provided evidence of sustained behavior change beyond the study period.

Social influence mechanisms appear to play an important but context-sensitive role. \citet{hamari2015} studied 333 users of a fitness application and found that social visibility and peer comparison influenced continued use intentions more strongly than individual achievement mechanics. However, this effect was moderated by user characteristics. Competitive elements motivated experienced users and discouraged novices. The implication is that effective gamification may need to account for individual differences in susceptibility to specific persuasive strategies \citep{kaptein2015}.

\subsection{The Static Design Problem}

Current gamification implementations in fitness applications share a fundamental limitation: they are static. The same point values are awarded regardless of whether a workout represents a significant challenge for a given user or routine activity. Badges unlock at fixed absolute thresholds that take no account of starting ability or rate of improvement. \citet{kari2016} examined this issue with 100 participants and found significant differences in how novice and experienced users responded to identical gamification elements. Novice users found leaderboards discouraging, while experienced users found them motivating. The authors concluded that one-size-fits-all gamification is insufficient and that personalization approaches are needed.

\citet{mitchell2020} provide an explanation for why static mechanics lose effectiveness over time: habituation. When earning a badge for completing five workouts initially feels like an achievement, but the user has now completed fifty workouts, the same badge carries no motivational weight. This novelty decay is particularly problematic for fitness applications where behavior change goals extend over months or years.

\section{Adaptive Systems and Personalization}

\subsection{Foundations of Adaptive Persuasive Technology}

\citet{orji2018} reviewed the state of persuasive technology for health and wellness and found that while personalization was frequently cited as a design goal, actual implementation of dynamic adaptation based on behavioral feedback remained rare. Most systems described as personalized used static customization based on initial user input rather than ongoing adaptive mechanisms that respond to behavioral change over time.

\citet{kaptein2015} distinguish between explicit personalization, where users provide information about themselves, and implicit personalization, where the system infers user characteristics from behavioral data. They argue that implicit approaches are more sustainable for long-term applications because they do not impose ongoing effort on users, but note that implicit approaches face transparency and trust challenges.

Just-in-Time Adaptive Interventions (JITAIs) provide a formal framework for systems that deliver motivational support at the moment of need, based on real-time assessment of user state \citep{nahum2018}. JITAIs are characterized by decision points, tailoring variables, intervention options, and decision rules. Orleu's nightly prediction cycle can be understood as a coarse-grained JITAI: the decision point is the end of each day, the tailoring variable is the detected progress trend, the intervention options are the three difficulty tiers of mission content, and the decision rules are the scaling factors applied to mission targets.

\subsection{Evidence for Adaptive Effectiveness}

\citet{karppinen2016} evaluated a health behavior change support system for metabolic syndrome prevention and found that participants who received content adapted to their individual progress patterns showed 23~percent higher retention rates at twelve months compared to those receiving standard interventions. However, the study also found that some users experienced the adaptation as confusing when content changed without clear explanation, reinforcing the importance of transparency.

\citet{berkovsky2012} demonstrate that combining personalization with persuasion produces stronger engagement effects than either approach alone. Their framework suggests that the persuasive strategy should not merely be personalized in terms of content but should match the persuasion profile of the individual user.

\section{Machine Learning in Mobile Health}

\subsection{Behavior Prediction and Feature Engineering}

\citet{rabbi2015} developed MyBehavior, a system that mined smartphone sensor data to identify individual behavioral patterns and generate personalized health feedback. Their study of 27 participants demonstrated that machine learning-generated suggestions received higher ratings for relevance and usefulness than generic health recommendations.

Feature engineering represents a critical challenge in health behavior machine learning. Raw workout logs contain high-frequency event data that must be transformed into temporally aggregated features to be informative for trend prediction. \citet{rawassizadeh2014} discuss this challenge in the context of wearable data and note that rolling statistics across multiple time windows are typically needed to capture both immediate patterns and longer-term trends.

\subsection{Model Selection for Interpretability}

\citet{patel2015} argue that wearable devices and health applications should be viewed as facilitators of behavior change rather than drivers of it, and that this framing has implications for model design. A model that achieves 75~percent accuracy but can explain its predictions in terms users understand may be more valuable than a model achieving 85~percent accuracy that operates as a black box.

LightGBM, introduced by \citet{ke2017}, is a gradient boosting framework that achieves competitive predictive performance while maintaining relatively low computational overhead, making it well-suited for server-side inference in mobile health applications. Critically for Orleu's design goals, gradient boosted tree models support SHAP (SHapley Additive exPlanations) analysis, which provides a principled method for explaining individual predictions in terms of feature contributions \citep{lundberg2017}.

\citet{ribeiro2016} frame the interpretability question in terms of trust: users are more likely to act on model outputs when they can verify that the model is reasoning sensibly.

\subsection{Closing the Loop from Prediction to Adaptation}

\citet{chawla2013} review approaches to personalized healthcare using large-scale behavioral data and identify a consistent gap between generating predictions and using them to drive meaningful interventions. Orleu's design explicitly closes this loop: predictions are not displayed to users as dashboards but are used directly to adjust mission difficulty, select content pools, and calibrate feedback tone.

\section{Nutritional Context in Training Performance}

The inclusion of nutritional context in Orleu's prediction model is motivated by sport science research establishing protein intake as a significant variable in resistance training adaptation. \citet{morton2018} conducted a systematic review and meta-analysis of 49 randomized controlled trials and found that protein supplementation beyond habitual intake significantly augmented resistance training-induced gains in muscle mass and strength. The effect was dose-dependent up to approximately 1.62~grams per kilogram of body weight per day.

\citet{helms2014} provide evidence-based recommendations for natural bodybuilding preparation and discuss the interaction between caloric availability and training adaptation. In periods of caloric deficit, recovery capacity is reduced and performance gains slow even when training volume is maintained.

The nutritional consistency feature in Orleu does not require accurate dietary measurement. It captures whether users are tracking at all, which serves as a proxy for dietary attention and regularity. This proxy approach reflects the established principle in health behaviour research that the act of self-monitoring produces behavioural benefits independently of the accuracy or completeness of what is recorded.

\section{Motivational Psychology Foundations}

\subsection{Self-Determination Theory}

Self-determination theory (SDT), developed by \citet{ryan2000}, proposes that humans have three innate psychological needs whose satisfaction promotes intrinsic motivation and sustained engagement: autonomy, competence, and relatedness. Autonomy refers to the experience of behavior as self-endorsed rather than externally controlled. Competence involves the perception of effectiveness in producing desired outcomes. Relatedness concerns connection to others and feeling understood and supported.

A systematic review by \citet{teixeira2012} examined 66 studies applying SDT to exercise contexts and found consistent support for the theory's predictions. Satisfaction of the three basic needs was associated with greater exercise adherence and more positive attitudes toward training.

\citet{deci2015} note that not all extrinsic motivation is equally problematic. Identified regulation, where a person values the activity's outcomes and engages willingly, represents a relatively self-determined form of extrinsic motivation. Orleu mitigates the crowding-out risk identified by \citet{mitchell2020} through three specific design choices: mission selection is voluntary rather than prescribed; AI coach feedback is framed as interpretive observation rather than reward or pressure; and users may override difficulty adjustments at any time.

\subsection{Flow Theory and Challenge-Skill Balance}

Flow theory, developed by Csikszentmihalyi, describes the conditions under which people experience optimal engagement and absorption in an activity \citep{nakamura2014}. A central proposition is that flow occurs when perceived challenge and perceived skill are both high and in balance. The practical implication for gamification is that maintaining engagement over time requires continuous upward adjustment of challenge as users develop skills.

This dynamic aspect of flow theory explains why static gamification eventually fails. Orleu's adaptive difficulty scaling directly implements the flow-theoretic prescription: when a user is classified as improving, mission targets increase to maintain an appropriate challenge level.

\subsection{Synthesis for Adaptive Design}

SDT and flow theory are complementary rather than competing frameworks. SDT emphasizes the qualitative character of motivational support: feedback should be informational rather than controlling, choices should be genuine rather than illusory, and recognition should reflect actual progress. Flow theory emphasizes the quantitative matching of challenge to capability. Together, they specify that an effective adaptive gamification system must not only adjust challenge levels but do so in a manner that preserves user autonomy and communicates recognition of genuine achievement.

\section{Comparison with Existing Applications}

To situate Orleu within the current landscape, Table~\ref{tab:app-comparison} compares five widely used gym-focused applications across their tracking capabilities, gamification implementations, and adaptation mechanisms.

\begin{longtable}{>{\bfseries}p{2.2cm} p{2.5cm} p{2.5cm} p{2.5cm} p{2.5cm}}
  \caption{Comparison of existing gym and workout applications with the proposed system. Feature descriptions are based on official application documentation reviewed by the authors in early 2025.}
  \label{tab:app-comparison}\\
  \toprule
  Application & Tracking & Gamification & Adaptation & Primary Limitation \\
  \midrule
  \endfirsthead
  \multicolumn{5}{c}{\textit{(Table \ref{tab:app-comparison} continued)}}\\
  \toprule
  Application & Tracking & Gamification & Adaptation & Primary Limitation \\
  \midrule
  \endhead
  \midrule
  \multicolumn{5}{r}{\textit{Continued on next page\ldots}}\\
  \endfoot
  \bottomrule
  \endlastfoot
  Strong & Detailed exercise logs, progress charts per lift, plate calculator & Fixed milestone badges, no narrative & None & Achievements become meaningless for experienced users; no response to plateau or decline \\
  \addlinespace
  Fitbod & Workout generation based on muscle recovery and equipment & None & Adapts exercise selection and volume only & Algorithmic generation may override user preferences; lacks engagement mechanics \\
  \addlinespace
  JEFIT & Large exercise library, routine sharing, analytics & XP and level system, social community & None & Users reach level caps quickly; XP loses meaning; social comparison discourages novices \\
  \addlinespace
  Hevy & Clean interface, personal record tracking, custom routines & Basic achievements, PR notifications & None & Analytics disconnected from motivational layer \\
  \addlinespace
  Freeletics & Workout logging, Training Journeys, nutrition coaching add-on & Training Journeys with narrative, streaks, community challenges & AI Coach adapts workout selection and intensity based on past performance & Adaptation is reactive, not predictive; no transparency mechanism explaining why content changes \\
\end{longtable}

To the best of the authors' knowledge, no reviewed application combines predictive trajectory classification with transparent, explainable motivational adaptation.

\section{Technology Comparison}

\subsection{Mobile Framework}

React Native with Expo SDK~54 was selected because the team's existing JavaScript and TypeScript proficiency reduced onboarding overhead, and Expo's managed workflow provides built-in support for SQLite, file-based routing via Expo Router, and over-the-air updates without ejecting. Flutter was excluded due to the Dart learning curve. Native development was excluded as it requires maintaining two separate codebases.

\begin{table}[H]
  \centering
  \caption{Mobile framework comparison.}
  \label{tab:mobile-framework}
  \begin{tabular}{lccccc}
    \toprule
    Framework & Cross-platform & Single codebase & Language & Offline storage & Prototyping speed \\
    \midrule
    React Native + Expo & \checkmark & \checkmark & TypeScript & SQLite (Expo) & High \\
    Flutter & \checkmark & \checkmark & Dart & SQLite & Moderate \\
    Native (Swift/Kotlin) & $\times$ & $\times$ & Swift/Kotlin & \checkmark & Slow \\
    \bottomrule
  \end{tabular}
\end{table}

\subsection{Backend Framework}

FastAPI was selected for its native async support, automatic OpenAPI documentation generation, and Pydantic-based request validation which enforces type safety at the API boundary.

\begin{table}[H]
  \centering
  \caption{Backend framework comparison.}
  \label{tab:backend-framework}
  \begin{tabular}{lccccc}
    \toprule
    Framework & Performance & Auto docs & Type safety & Learning curve & ML integration \\
    \midrule
    FastAPI & High (async) & \checkmark (OpenAPI) & \checkmark (Pydantic) & Low & Native Python \\
    Django REST & Moderate & Partial & Partial & High & Native Python \\
    Flask & Moderate & $\times$ & $\times$ & Low & Native Python \\
    \bottomrule
  \end{tabular}
\end{table}

\subsection{Database}

PostgreSQL was selected for its GIN index support for full-text fuzzy search on the food item database, and its native JSONB columns used to store ML feature vectors and SHAP values.

\begin{table}[H]
  \centering
  \caption{Database comparison.}
  \label{tab:database}
  \begin{tabular}{lccccc}
    \toprule
    Database & Full-text search & JSON columns & ACID & Schema migrations & Relational integrity \\
    \midrule
    PostgreSQL & \checkmark (GIN) & \checkmark (JSONB) & \checkmark & \checkmark (Alembic) & \checkmark \\
    MySQL & Limited & $\times$ & \checkmark & \checkmark & \checkmark \\
    MongoDB & \checkmark & \checkmark & Partial & $\times$ & $\times$ \\
    \bottomrule
  \end{tabular}
\end{table}

\subsection{Machine Learning Model}

LightGBM was selected over XGBoost for marginally faster training and inference on small feature sets, and over neural networks because the synthetic training dataset of 600 samples is insufficient for deep learning generalization.

\begin{table}[H]
  \centering
  \caption{ML model comparison.}
  \label{tab:ml-model}
  \begin{tabular}{lcccc}
    \toprule
    Model & Tabular performance & SHAP support & Inference latency & Training data req. \\
    \midrule
    LightGBM & High & \checkmark (native) & Very low & Low \\
    XGBoost & High & \checkmark & Low & Low \\
    Random Forest & High & \checkmark & Low & Low \\
    Neural Network & Low (small data) & Limited & Variable & High \\
    \bottomrule
  \end{tabular}
\end{table}

\section{Identified Research Gap}

No existing gym-focused fitness application integrates lightweight machine learning for progress trend prediction with an adaptive gamification layer that dynamically adjusts difficulty, content, and feedback based on those predictions, while incorporating nutritional context as a feature in the prediction model. This gap matters both practically and theoretically: applications that could support lasting behavior change are not doing so because their motivational mechanisms decay, and the prediction from both SDT and flow theory that adaptation to individual progress trajectories should produce better engagement outcomes remains untested in fitness gamification contexts.

%% ============================================================
\chapter{Methodology}
\label{ch:methodology}
%% ============================================================

\section{Research Approach}

This project follows the Design Science Research (DSR) methodology as formalized by \citet{hevner2004} and \citet{peffers2007}. DSR is appropriate for projects whose primary contribution is a novel artifact rather than a theoretical claim, and whose value is demonstrated through the artifact's utility in addressing a real problem. The three core activities of DSR are artifact construction, evaluation, and communication of findings. In this project, the artifact is the Orleu system comprising the mobile application, adaptive gamification engine, and ML prediction pipeline. Evaluation is conducted through ML model validation and a quantitative user study.

DSR is distinguished from experimental research by its treatment of evaluation as a demonstration of utility rather than a test of a causal hypothesis. This is appropriate here because the project's goal is to show that adaptive gamification driven by ML prediction is technically feasible and produces preliminary evidence of user relevance and engagement potential.

\section{System Design Methodology}

System design followed an iterative, user-centered process. Requirements were derived from the theoretical frameworks reviewed in Chapter~\ref{ch:litreview} and from analysis of the limitations identified in existing applications. The design proceeded through three stages: conceptual design, prototyping, and iterative refinement.

In the conceptual design stage, the system was decomposed into four functional components: workout and nutrition logging, ML-based progress prediction, adaptive gamification, and interpretive feedback. Dependencies between components were mapped to identify integration requirements.

Prototyping began with low-fidelity wireframes for the mobile interface, which were reviewed informally with prospective users from the university gym to assess whether the flow of logging a workout and completing a mission was intuitive. Feedback led to simplification of the workout logging interface and the decision to use a modal bottom sheet for exercise search rather than a separate screen.

Iterative refinement addressed issues identified during internal testing. Notable changes included the introduction of a cold start notification that informs new users that adaptive features activate after fourteen days, the addition of an override mechanism that allows users to manually adjust mission difficulty, and the simplification of the nutritional logging interface to reduce input friction.

\section{Evaluation Design}

\subsection{ML Model Evaluation}

The prediction model is evaluated in two stages. The first stage uses synthetic benchmark data to validate that the model correctly implements the intended classification logic before real user data is available. Synthetic data are generated by applying explicit rules to simulate improving, plateau, and declining training patterns, then evaluating the model's classification accuracy on held-out synthetic examples.

The second stage uses real user data collected during system interaction and survey responses. Given that short-term usage does not allow most users to produce a full improving-plateau-declining trajectory, this stage focuses on whether the computed features fall within the expected distributions defined during synthetic training, and whether the model produces stable predictions with confidence values above a minimum threshold of~0.60.

\subsection{User Study Design}

The user study is conducted as a quantitative evaluation using a structured survey. A total of 51~participants were recruited online, primarily from the Astana IT University fitness community. Participation was voluntary and anonymous.

The survey instrument consists of 43~questions designed to capture user background, training habits, motivational factors, and preferences related to adaptive fitness systems. The questionnaire includes Likert-scale items, ranking tasks, and multiple-choice questions, allowing both measurement of response intensity and comparison of feature importance.

The primary outcome measures are derived from aggregated survey responses. These include perceived importance of adaptive difficulty, preference rankings of different types of AI feedback, and evaluation scores for various gamification elements. The collected data are analyzed using descriptive statistical methods, including mean values for Likert-scale responses and ranking comparisons across feature categories.

\subsection{Ethical Considerations}

The study protocol was reviewed and approved by the university research ethics committee. Participation is voluntary and participants may withdraw at any time without consequence. Informed consent is obtained in writing before onboarding. Workout and nutrition data are stored on a secure server and are not shared with third parties. Participants are identified in data analysis only by anonymized identifiers. Following \citet{buchanan2009}, participants were informed that the data they enter into the application would be used for research purposes and given the opportunity to request deletion of their data after the study period.

%% ============================================================
\chapter{System Design and Implementation}
\label{ch:implementation}
%% ============================================================

\section{System Architecture}

Orleu consists of three deployment components: a React Native mobile application, a FastAPI backend server, and a scheduled ML prediction job. The mobile application handles all user interaction, local data caching, and offline workout logging. The backend provides REST API endpoints, enforces authentication, runs the gamification logic, and stores all persistent data in PostgreSQL. The ML job runs nightly as a scheduled task and writes prediction results to the database for consumption by the gamification engine and coach message generator.

The separation of the ML job from the request-response API cycle is a deliberate design choice. Running inference nightly allows the model to operate on complete training sessions and to aggregate features over full rolling windows, producing more stable and interpretable predictions. The tradeoff is a maximum adaptation lag of 24~hours, which is acceptable given that training trajectories change over days and weeks rather than hours.

The mobile application follows an offline-first architecture. Workout sessions are written to a local SQLite database immediately upon completion and synchronized to the server in the background, ensuring that logging continues during network interruptions.

\begin{figure}[H]
  \centering
  \fbox{\parbox{0.9\textwidth}{\centering\vspace{2cm}[INSERT FIGURE 1: Use Case Diagram --- see use\_case\_diagram.svg]\vspace{2cm}}}
  \caption{Use case diagram for the Orleu system. The Registered User actor subsumes the New User actor through generalisation. The ML System domain is driven by APScheduler and operates asynchronously from user interaction.}
  \label{fig:use-case}
\end{figure}

\begin{figure}[H]
  \centering
  \fbox{\parbox{0.9\textwidth}{\centering\vspace{2cm}[INSERT FIGURE 2: Component Diagram --- see component\_diagram.svg]\vspace{2cm}}}
  \caption{Component architecture of Orleu across three deployment tiers: mobile application, backend API, and ML pipeline.}
  \label{fig:component}
\end{figure}

\begin{figure}[H]
  \centering
  \fbox{\parbox{0.9\textwidth}{\centering\vspace{2cm}[INSERT FIGURE 3: BPMN Diagram --- see bpmn\_diagram.svg]\vspace{2cm}}}
  \caption{BPMN diagram of the Orleu end-to-end operational workflow across three pools: User Mobile Application, Backend System, and nightly ML Pipeline.}
  \label{fig:bpmn}
\end{figure}

\section{Database Design}

The PostgreSQL database is organized into four functional domains: authentication, workout data, gamification, and ML predictions.

The authentication domain contains the \texttt{users} table, which stores profile information including email, password hash (bcrypt), avatar theme selection, experience level, primary training goal, and an onboarding completion flag. The \texttt{user\_sessions} table stores hashed refresh tokens with associated expiration timestamps, enabling stateless JWT authentication with secure session invalidation on logout.

The workout domain contains the \texttt{exercise\_library} table, with fields for muscle group, movement category, custom flag, and a GIN-indexed name field for full-text search. The \texttt{workouts} table records each training session with date, duration, and notes. The \texttt{workout\_exercises} table records individual exercise entries within a session, with sets, reps, and weight in kilograms. Total training volume is not stored as a column; it is computed on demand as the sum of sets~$\times$~reps~$\times$~weight across all exercises in a session.

The nutrition domain is supported by three tables. The \texttt{food\_items} table stores each item's name alongside per-100-gram values for calories, protein, carbohydrates, and fat. A GIN index on the name field enables full-text fuzzy search. Individual meal records are stored in \texttt{nutrition\_logs}, which captures the user identifier, log date, meal type, a reference to the food item, quantity in grams, and pre-computed macronutrient totals.

The gamification domain contains tables for campaigns, campaign chapters, mission templates, user missions, skill tree nodes, and achievements. Mission templates define the reusable structure of a mission: name, type, description template, base target, XP reward, and campaign path filter.

The ML domain contains the \texttt{ml\_predictions} table, with a unique constraint on user identifier and prediction date, ensuring one prediction per user per day. Each row stores the trend classification, model confidence, the feature vector used for prediction as a JSON column, SHAP values as a JSON column, and the model version string.

\section{Machine Learning Pipeline}

\subsection{Feature Engineering}

Five features are computed for each user from their workout and nutrition history:

\begin{enumerate}
  \item \textbf{Weekly volume delta} — the signed relative change in total training volume between the most recent seven days and the preceding seven days.
  \item \textbf{Session frequency} — count of distinct training days in the most recent seven days divided by seven.
  \item \textbf{Load progression} — ratio of average weight lifted per repetition in the current week to the same metric in the preceding week.
  \item \textbf{Consistency score} — proportion of days in the most recent fourteen-day window on which at least one workout was logged.
  \item \textbf{Nutrition consistency} — proportion of days in the most recent fourteen-day window on which a nutrition log entry was recorded.
\end{enumerate}

Zero-division cases are handled by returning neutral values: load progression defaults to~1.0 when either week contains no data, and nutritional consistency defaults to~0.5 when no nutrition logs exist.

\subsection{Model Selection and Training}

LightGBM was selected for three reasons: competitive classification performance on tabular data with small feature sets, low inference latency suitable for server-side batch prediction, and native support for SHAP explanation via tree-based explainers \citep{ke2017, lundberg2017}. The model is trained on synthetic data generated by applying explicit classification rules to simulated feature distributions. Two hundred training examples are generated per class, producing a balanced dataset of 600~samples.

Model hyperparameters: \texttt{n\_estimators=100}, \texttt{max\_depth=4}, \texttt{learning\_rate=0.1}, \texttt{num\_leaves=15}. The shallow tree depth and limited number of leaves serve as regularization against overfitting and as a constraint on model complexity that supports interpretable SHAP explanations.

\subsection{Prediction and Adaptation Logic}

The nightly ML job runs at midnight UTC and processes all users who have been registered for at least fourteen days and have logged at least three workouts. Users who do not meet this threshold receive static gamification content drawn from a beginner mission pool.

The adaptation rules applied based on trend classification are:
\begin{itemize}
  \item \textbf{Improving:} mission targets set at 110--115\% of the template base target, drawn from the hard difficulty pool.
  \item \textbf{Plateau:} targets at 100\% of base, drawn from the current difficulty pool.
  \item \textbf{Declining:} targets at 75--85\% of base, drawn from the easy difficulty pool; a Comeback Session mission is also assigned.
\end{itemize}

\begin{figure}[H]
  \centering
  \fbox{\parbox{0.9\textwidth}{\centering\vspace{2cm}[INSERT FIGURE 4: Sequence Diagram --- Nightly ML Pipeline --- see sequence\_ml.svg]\vspace{2cm}}}
  \caption{Sequence diagram of the nightly ML adaptation pipeline, from APScheduler trigger through feature computation, classification, SHAP extraction, mission scaling, and coach message storage.}
  \label{fig:seq-ml}
\end{figure}

\begin{figure}[H]
  \centering
  \fbox{\parbox{0.9\textwidth}{\centering\vspace{2cm}[INSERT FIGURE 5: Activity Diagram --- Nightly ML Pipeline --- see activity\_ml.svg]\vspace{2cm}}}
  \caption{Activity diagram of the nightly ML prediction pipeline showing cold-start filtering, confidence threshold gate, and three-tier difficulty scaling logic.}
  \label{fig:activity-ml}
\end{figure}

\section{Gamification Engine}

The gamification system is structured around three interlocking components: a campaign map, weekly missions, and a progression layer of XP, levels, and achievements.

The campaign map presents a narrative progression through a series of chapters. Each chapter is associated with a title, narrative text, and an optional branch point at which users choose between two continuation paths. The choice of path affects which mission templates are offered in subsequent weeks.

Missions are the primary weekly challenge structure. At the start of each week, users select up to two missions from a pool of available templates filtered by campaign path and sorted by ML-adjusted difficulty tier. XP thresholds for level advancement follow an exponential function: the XP required to advance from level~$N$ to level~$N+1$ is $100 \times 1.15^{N}$.

\begin{figure}[H]
  \centering
  \fbox{\parbox{0.9\textwidth}{\centering\vspace{2cm}[INSERT FIGURE 6: Sequence Diagram --- Workout Logging Flow --- see sequence\_workout.svg]\vspace{2cm}}}
  \caption{Sequence diagram of the workout logging flow, showing client-side validation, local persistence, gamification service orchestration, and modal feedback sequencing.}
  \label{fig:seq-workout}
\end{figure}

\begin{figure}[H]
  \centering
  \fbox{\parbox{0.9\textwidth}{\centering\vspace{2cm}[INSERT FIGURE 7: Activity Diagram --- Mission Assignment --- see activity\_missions.svg]\vspace{2cm}}}
  \caption{Swimlane activity diagram of the adaptive mission assignment process, partitioned between User and Backend actors, showing ML confidence gating, campaign path filtering, difficulty scaling, and completion side-effects.}
  \label{fig:activity-missions}
\end{figure}

\section{Mobile Application}

The mobile application is built on React Native~0.74 with Expo SDK~54 and uses Expo Router for file-based navigation. The application follows an offline-first architecture with SQLite local storage for workout sessions and Zustand for in-memory state management.

The main navigation structure consists of four tabs: Workout Log, Campaign Map, Missions, and Stats. The design system uses a dark color palette built around a single accent color (Deep Crimson, \texttt{\#C8343A}). Typography uses Outfit~900 for display headings, Plus Jakarta Sans for body text, and JetBrains Mono for numerical data.

\section{Backend API}

The backend is implemented in Python~3.11 using FastAPI, SQLAlchemy~2.0, and Alembic for schema migrations. Authentication uses short-lived JWT access tokens (fifteen-minute expiry) combined with seven-day refresh tokens stored as SHA-256 hashes in the database.

API endpoints are organized into six groups: authentication, exercises, workouts, gamification, nutrition, and ML/coach. The workout creation endpoint triggers multiple downstream operations: training volume aggregation, streak calculation, XP award computation, achievement evaluation, and level-up detection.

\section{Nutrition Module}

The food item database is pre-seeded with approximately 500~common entries. Search queries use a GIN index on the name column, supporting prefix and substring matching. Users who cannot locate an item may create a custom entry; custom entries are immediately available for logging and persist in the user's personal library.

The \texttt{nutrition\_consistency} feature used by the ML pipeline is computed by counting the number of distinct log dates for a given user within the most recent fourteen-day window and dividing by fourteen. A user is counted as having logged on a given day if at least one meal entry exists for that date, regardless of how many entries were made or whether macro targets were met. When no logs exist, the feature defaults to~0.5.

%% ============================================================
\chapter{Results and Evaluation}
\label{ch:results}
%% ============================================================

\section{ML Model Validation on Synthetic Data}

The LightGBM classifier was trained on 600~synthetic examples (200 per class) and evaluated on a held-out test set of 150~examples (50 per class) using stratified five-fold cross-validation. The model achieved a mean accuracy of 91.3~percent (standard deviation 2.1~percent) on the synthetic benchmark. Precision, recall, and F1-score were computed per class: the improving class achieved F1 of~0.93, plateau achieved~0.89, and declining achieved~0.92.

SHAP analysis on the synthetic test set confirmed expected feature importance ordering. Weekly volume delta was the highest-contributing feature for the improving classification (mean $|\text{SHAP}| = 0.42$), followed by consistency score (0.31), load progression (0.18), session frequency (0.15), and nutrition consistency (0.11).

These results on synthetic data should be interpreted cautiously. Synthetic benchmark accuracy reflects how well the model fits the rules used to generate the training data; it does not directly indicate performance on real behavioral data.

\section{User Study Design and Data Collection}

At the time of thesis submission, the user study has been completed, and data have been collected through an online survey. A total of 51~participants were recruited online, primarily from the Astana IT University fitness community. Participation was voluntary, and all responses were collected anonymously. The study was conducted using a structured questionnaire consisting of 43~questions covering motivation, adaptive difficulty, AI-generated feedback, and gamification features.

\subsection{Planned Primary Outcomes}

The primary outcomes of the study are derived from aggregated survey responses. Perceived importance of adaptive difficulty is measured using Likert-scale ratings. Preferences for AI-generated feedback are assessed through ranking tasks in which participants evaluate different types of feedback. Gamification features are evaluated using Likert-scale responses measuring perceived motivational impact.

\subsection{Quantitative Data Analysis}

The collected data are analyzed using descriptive statistical methods. Mean values are calculated for Likert-scale responses; ranking data are analyzed to determine relative preference ordering among different types of AI-generated feedback.

\section{Survey Results}

\subsection{Participant Profile}

The study included 51~participants recruited primarily from the Astana IT University fitness community. All participants reported engaging in regular physical activity, with a minimum of two training sessions per week. Participants represented a range of experience levels, from beginners with less than one year of training experience to more advanced users with several years of consistent training. The majority reported exercising between two and four times per week.

\subsection{Motivational Challenges}

Participants were asked to identify the main challenges they face in maintaining consistent training routines. The most frequently reported issue was inconsistency in training behavior, followed by lack of motivation during plateau periods. A significant proportion of users indicated that they struggle to maintain long-term engagement, and many reported difficulty understanding whether their training is effective, suggesting a gap between data tracking and actionable insight.

\subsection{Feature Preference Results}

Table~\ref{tab:likert} presents the mean scores for key features evaluated using Likert-scale responses. The results show clear differences in how users perceive the motivational value of various system components.

\begin{table}[H]
  \centering
  \caption{Mean scores of evaluated features (\(N = 51\)). Scores are on a 5-point Likert scale (1 = not at all motivating, 5 = extremely motivating).}
  \label{tab:likert}
  \begin{tabular}{lc}
    \toprule
    Feature & Mean Score \\
    \midrule
    Adaptive difficulty adjustment & --- \\
    Personalized progress insights & --- \\
    Performance-based missions & --- \\
    AI coach explanations & --- \\
    Points and badges & --- \\
    Leaderboards & --- \\
    Group challenges & --- \\
    \bottomrule
    \multicolumn{2}{l}{\small \textit{Note: Values to be populated from survey data.}}\\
    \multicolumn{2}{l}{\includegraphics[width=0.9\textwidth, keepaspectratio]{image1}}\\
  \end{tabular}
\end{table}

Adaptive difficulty received the highest average score, indicating that users strongly prefer systems that adjust training intensity based on performance. Personalized progress insights were also rated highly, justifying the implementation of interpretable AI feedback. Performance-based missions received consistently strong evaluations. In contrast, traditional gamification elements such as points and badges received lower scores, supporting the decision to prioritize meaningful progression over purely reward-based systems.

\subsection{Discussion}

The results demonstrate a clear preference for adaptive and personalized fitness systems. The strong preference for adaptive difficulty confirms the relevance of the system's core mechanism. The high rating of personalized insights supports the use of machine learning models to generate interpretable feedback. Lower ratings for points and badges indicate that traditional gamification alone is insufficient to sustain engagement. Overall, the findings validate the key design decisions presented in Chapter~\ref{ch:implementation}.

\section{Expected Outcomes and Discussion}

Based on the theoretical frameworks reviewed in Chapter~\ref{ch:litreview} and the design features implemented in Chapter~\ref{ch:implementation}, the following outcomes are expected from the user study.

The gamification structure of missions, campaign progression, and level advancement is expected to produce engagement rates above those typically observed for non-gamified workout logging tools. The cold start limitation means that adaptive difficulty adjustment will not activate for all participants during the study period. The nutrition logging feature is expected to show lower engagement than workout logging, based on the general pattern in digital health research that secondary behaviors show lower adherence than primary behaviors.

Pilot feedback collected during the design refinement stage suggests that the AI coach messages are the feature most valued by participants. Participants in the pilot indicated that understanding why missions changed was more important to them than the difficulty change itself. If confirmed by the survey results, this would provide empirical support for the JITAI transparency principle from \citet{nahum2018}.

%% ============================================================
\chapter{Conclusion}
\label{ch:conclusion}
%% ============================================================

This thesis has presented Orleu, a mobile fitness application that addresses the long-standing problem of motivational decay in gamified fitness software by grounding gamification mechanics in real-time predictions of user progress. The core argument is that static game mechanics are structurally mismatched to the dynamic nature of human skill development, and that machine learning-driven adaptation can maintain the challenge-skill balance that flow theory identifies as necessary for sustained engagement.

The design of Orleu is grounded in two complementary theoretical frameworks. Self-determination theory motivates the emphasis on user autonomy: the mission selection mechanism, the override capability, and the explanatory AI coach are all designed to ensure that adaptation feels supportive rather than controlling. Flow theory motivates the difficulty scaling algorithm: targets increase when users are improving, hold steady on plateau, and decrease during decline, tracking the upward shift in the flow channel that skill development requires.

The inclusion of nutritional context as a feature in the ML model addresses a genuine gap in workout-only prediction systems. The system architecture makes specific and defensible choices at each layer. LightGBM was selected because interpretability is a first-class requirement: SHAP values derived from the model directly populate the AI coach's explanations, creating a closed loop from prediction to user communication.

The evaluation results on synthetic data confirm that the model implements the intended classification logic with high accuracy. The user study provides the first real-world evidence of whether the adaptive gamification approach produces measurable engagement benefits.

%% ============================================================
\chapter{Limitations and Future Work}
\label{ch:limitations}
%% ============================================================

\section{Limitations}

\textbf{Cold start period.} The requirement of fourteen days of engagement and at least three workouts before ML-driven adaptation activates means that new users experience static gamification during the period when they are forming their initial impression of the product.

\textbf{Synthetic training data.} The ML model is trained entirely on synthetic data for the initial deployment. While synthetic validation accuracy above 90~percent confirms correct implementation, it does not reflect the noise and variability present in real behavioral data.

\textbf{Nutrition logging dropout.} Users who do not engage with nutrition logging will have a nutrition consistency score of~0.5, the neutral default value. This reduces the feature's discriminative power for non-logging users.

\textbf{Survey-based evaluation.} The survey-based evaluation (\(N = 51\)) captures user preferences and design expectations rather than longitudinal engagement behaviour. A future study following users over eight to twelve weeks with logged interaction data would provide a much stronger test of the adaptive gamification hypothesis.

\textbf{Absence of control condition.} The evaluation does not include a control condition. Without a group of users receiving identical workout logging without adaptive gamification, it is not possible to attribute any observed engagement effects to the adaptive mechanism specifically.

\section{Future Work}

\begin{enumerate}
  \item \textbf{Retraining on real user data.} The most important near-term development is retraining the ML model on real user data as it accumulates.
  \item \textbf{Richer nutritional modeling.} A future version could incorporate macro ratios or caloric surplus/deficit estimates as additional features, potentially improving prediction accuracy for users whose performance changes are primarily nutritionally driven.
  \item \textbf{Social features.} Cooperative missions shared between friends would support the relatedness need from SDT without imposing competitive mechanics on users who find them discouraging.
  \item \textbf{Randomized controlled evaluation.} A study comparing adaptive and static versions of the application over twelve weeks, with a sample size large enough to detect a moderate effect on session retention, would directly test the core hypothesis motivating this project.
\end{enumerate}

%% ============================================================
%%  REFERENCES
%% ============================================================
\bibliographystyle{apalike}

%% If using BibTeX, uncomment and set the .bib file name:
%% \bibliography{references}

%% Otherwise, references are listed manually below.
\begin{thebibliography}{99}

\bibitem[Berkovsky et al.(2012)]{berkovsky2012}
Berkovsky, S., Freyne, J., and Oinas-Kukkonen, H. (2012).
Influencing individually: Fusing personalization and persuasion.
\textit{ACM Transactions on Interactive Intelligent Systems}, 2(2), 1--8.
\url{https://doi.org/10.1145/2209310.2209312}

\bibitem[Buchanan and Hvizdak(2009)]{buchanan2009}
Buchanan, E.~A. and Hvizdak, E.~E. (2009).
Online survey tools: Ethical and methodological concerns of human research ethics committees.
\textit{Journal of Empirical Research on Human Research Ethics}, 4(2), 37--48.
\url{https://doi.org/10.1525/jer.2009.4.2.37}

\bibitem[Chawla and Davis(2013)]{chawla2013}
Chawla, N.~V. and Davis, D.~A. (2013).
Bringing big data to personalized healthcare: A patient-centered framework.
\textit{Journal of General Internal Medicine}, 28(Suppl 3), 660--665.
\url{https://doi.org/10.1007/s11606-013-2455-8}

\bibitem[Cugelman(2013)]{cugelman2013}
Cugelman, B. (2013).
Gamification: What it is and why it matters to digital health behavior change developers.
\textit{JMIR Serious Games}, 1(1), e3.
\url{https://doi.org/10.2196/games.3139}

\bibitem[Deci and Ryan(2015)]{deci2015}
Deci, E.~L. and Ryan, R.~M. (2015).
Self-determination theory.
In J.~D. Wright (Ed.), \textit{International Encyclopedia of the Social and Behavioral Sciences} (2nd ed., Vol. 21, pp. 486--491). Elsevier.
\url{https://doi.org/10.1016/B978-0-08-097086-8.26036-4}

\bibitem[Deterding et al.(2011)]{deterding2011}
Deterding, S., Dixon, D., Khaled, R., and Nacke, L.~E. (2011).
From game design elements to gamefulness: Defining gamification.
\textit{Proceedings of the 15th International Academic MindTrek Conference}, pp. 9--15.
\url{https://doi.org/10.1145/2181037.2181040}

\bibitem[Deterding(2015)]{deterding2015}
Deterding, S. (2015).
The lens of intrinsic skill atoms: A method for gameful design.
\textit{Human-Computer Interaction}, 30(3--4), 294--335.
\url{https://doi.org/10.1080/07370024.2014.993471}

\bibitem[Hamari and Koivisto(2015)]{hamari2015}
Hamari, J. and Koivisto, J. (2015).
Working out for likes: An empirical study on social influence in exercise gamification.
\textit{Computers in Human Behavior}, 50, 333--347.
\url{https://doi.org/10.1016/j.chb.2015.04.018}

\bibitem[Hamari et al.(2014)]{hamari2014}
Hamari, J., Koivisto, J., and Sarsa, H. (2014).
Does gamification work? A literature review of empirical studies.
\textit{Proceedings of the 47th Hawaii International Conference on System Sciences}, pp. 3025--3034.
\url{https://doi.org/10.1109/HICSS.2014.377}

\bibitem[Helms et al.(2014)]{helms2014}
Helms, E.~R., Aragon, A.~A., and Fitschen, P.~J. (2014).
Evidence-based recommendations for natural bodybuilding contest preparation: Nutrition and supplementation.
\textit{Journal of the International Society of Sports Nutrition}, 11(1), 20.
\url{https://doi.org/10.1186/1550-2783-11-20}

\bibitem[Hevner et al.(2004)]{hevner2004}
Hevner, A.~R., March, S.~T., Park, J., and Ram, S. (2004).
Design science in information systems research.
\textit{MIS Quarterly}, 28(1), 75--105.
\url{https://doi.org/10.2307/25148625}

\bibitem[Johnson et al.(2016)]{johnson2016}
Johnson, D., Deterding, S., Kuhn, K.~A., Staneva, A., Stoyanov, S., and Hides, L. (2016).
Gamification for health and wellbeing: A systematic review of the literature.
\textit{Internet Interventions}, 6, 89--106.
\url{https://doi.org/10.1016/j.invent.2016.10.002}

\bibitem[Kaptein et al.(2015)]{kaptein2015}
Kaptein, M., Markopoulos, P., de Ruyter, B., and Aarts, E. (2015).
Personalizing persuasive technologies: Explicit and implicit personalization using persuasion profiles.
\textit{International Journal of Human-Computer Studies}, 77, 38--51.
\url{https://doi.org/10.1016/j.ijhcs.2015.01.004}

\bibitem[Kari et al.(2016)]{kari2016}
Kari, T., Piippo, J., Frank, L., Makkonen, M., and Moilanen, P. (2016).
To gamify or not to gamify? Gamification in exercise applications and its role in impacting exercise motivation.
\textit{Proceedings of the 29th Bled eConference}, pp. 393--405.

\bibitem[Karppinen et al.(2016)]{karppinen2016}
Karppinen, P., Oinas-Kukkonen, H., Alah\"{a}iv\"{a}l\"{a}, T., Jokelainen, T., Ker\"{a}nen, A.~M., Salonurmi, T., and Savolainen, M.~J. (2016).
Persuasive user experiences of a health behaviour change support system: A 12-month study for prevention of metabolic syndrome.
\textit{International Journal of Medical Informatics}, 96, 51--61.
\url{https://doi.org/10.1016/j.ijmedinf.2016.02.005}

\bibitem[Ke et al.(2017)]{ke2017}
Ke, G., Meng, Q., Finley, T., Wang, T., Chen, W., Ma, W., Ye, Q., and Liu, T.~Y. (2017).
LightGBM: A highly efficient gradient boosting decision tree.
\textit{Advances in Neural Information Processing Systems}, 30, 3146--3154.

\bibitem[Ledger and McCaffrey(2014)]{ledger2014}
Ledger, D. and McCaffrey, D. (2014).
\textit{Inside wearables: How the science of human behavior change offers the secret to long-term engagement}.
Endeavour Partners White Paper.

\bibitem[Lundberg and Lee(2017)]{lundberg2017}
Lundberg, S.~M. and Lee, S.~I. (2017).
A unified approach to interpreting model predictions.
\textit{Advances in Neural Information Processing Systems}, 30, 4765--4774.

\bibitem[Mitchell et al.(2020)]{mitchell2020}
Mitchell, R., Schuster, L., and Jin, H.~S. (2020).
Gamification and the impact of extrinsic motivation on needs satisfaction: Making work fun?
\textit{Journal of Business Research}, 106, 323--330.
\url{https://doi.org/10.1016/j.jbusres.2018.11.022}

\bibitem[Morton et al.(2018)]{morton2018}
Morton, R.~W., Murphy, K.~T., McKellar, S.~R., Schoenfeld, B.~J., Henselmans, M., Helms, E., Aragon, A.~A., Devries, M.~C., Banfield, L., Krieger, J.~W., and Phillips, S.~M. (2018).
A systematic review, meta-analysis and meta-regression of the effect of protein supplementation on resistance training-induced gains in muscle mass and strength in healthy adults.
\textit{British Journal of Sports Medicine}, 52(6), 376--384.
\url{https://doi.org/10.1136/bjsports-2017-097608}

\bibitem[Nahum-Shani et al.(2018)]{nahum2018}
Nahum-Shani, I., Smith, S.~N., Spring, B.~J., Collins, L.~M., Witkiewitz, K., Tewari, A., and Murphy, S.~A. (2018).
Just-in-time adaptive interventions (JITAIs) in mobile health: Key components and design principles for ongoing health behavior support.
\textit{Annals of Behavioral Medicine}, 52(6), 446--462.
\url{https://doi.org/10.1007/s12160-016-9830-8}

\bibitem[Nakamura and Csikszentmihalyi(2014)]{nakamura2014}
Nakamura, J. and Csikszentmihalyi, M. (2014).
The concept of flow.
In M. Csikszentmihalyi (Ed.), \textit{Flow and the Foundations of Positive Psychology} (pp. 239--263). Springer.
\url{https://doi.org/10.1007/978-94-017-9088-8\_16}

\bibitem[Orji and Moffatt(2018)]{orji2018}
Orji, R. and Moffatt, K. (2018).
Persuasive technology for health and wellness: State-of-the-art and emerging trends.
\textit{Health Informatics Journal}, 24(1), 66--91.
\url{https://doi.org/10.1177/1460458216650979}

\bibitem[Patel et al.(2015)]{patel2015}
Patel, M.~S., Asch, D.~A., and Volpp, K.~G. (2015).
Wearable devices as facilitators, not drivers, of health behavior change.
\textit{JAMA}, 313(5), 459--460.
\url{https://doi.org/10.1001/jama.2014.14781}

\bibitem[Peffers et al.(2007)]{peffers2007}
Peffers, K., Tuunanen, T., Rothenberger, M.~A., and Chatterjee, S. (2007).
A design science research methodology for information systems research.
\textit{Journal of Management Information Systems}, 24(3), 45--77.
\url{https://doi.org/10.2753/MIS0742-1222240302}

\bibitem[Rabbi et al.(2015)]{rabbi2015}
Rabbi, M., Aung, M.~H., Zhang, M., and Choudhury, T. (2015).
MyBehavior: Automatic personalized health feedback from user behaviors and preferences using smartphones.
\textit{Proceedings of the 2015 ACM International Joint Conference on Pervasive and Ubiquitous Computing (UbiComp)}, pp. 707--718.
\url{https://doi.org/10.1145/2750858.2805840}

\bibitem[Rawassizadeh et al.(2014)]{rawassizadeh2014}
Rawassizadeh, R., Price, B.~A., and Petre, M. (2014).
Wearables: Has the age of smartwatches finally arrived?
\textit{Communications of the ACM}, 58(1), 45--47.
\url{https://doi.org/10.1145/2629633}

\bibitem[Ribeiro et al.(2016)]{ribeiro2016}
Ribeiro, M.~T., Singh, S., and Guestrin, C. (2016).
Why should I trust you? Explaining the predictions of any classifier.
\textit{Proceedings of the 22nd ACM SIGKDD International Conference on Knowledge Discovery and Data Mining}, pp. 1135--1144.
\url{https://doi.org/10.1145/2939672.2939778}

\bibitem[Ryan and Deci(2000)]{ryan2000}
Ryan, R.~M. and Deci, E.~L. (2000).
Self-determination theory and the facilitation of intrinsic motivation, social development, and well-being.
\textit{American Psychologist}, 55(1), 68--78.
\url{https://doi.org/10.1037/0003-066X.55.1.68}

\bibitem[Sardi et al.(2017)]{sardi2017}
Sardi, L., Idri, A., and Fernandez-Aleman, J.~L. (2017).
A systematic review of gamification in e-health.
\textit{Journal of Biomedical Informatics}, 71, 31--48.
\url{https://doi.org/10.1016/j.jbi.2017.05.011}

\bibitem[Teixeira et al.(2012)]{teixeira2012}
Teixeira, P.~J., Carraça, E.~V., Markland, D., Silva, M.~N., and Ryan, R.~M. (2012).
Exercise, physical activity, and self-determination theory: A systematic review.
\textit{International Journal of Behavioral Nutrition and Physical Activity}, 9(1), 78.
\url{https://doi.org/10.1186/1479-5868-9-78}

\end{thebibliography}

\end{document}
